
\documentclass{article}


\usepackage[style=alphabetic]{biblatex}
\addbibresource{main.bib}


\usepackage{listings}
\lstdefinestyle{lst-default}{ 
    numberstyle=\tiny,
    basicstyle=\ttfamily\footnotesize,     
    keepspaces=true,                 
    numbers=left,                    
    numbersep=10pt,                  
    showspaces=false,                
    showstringspaces=false,
    showtabs=false,                  
    tabsize=2,
    xleftmargin=0.25in,
    xrightmargin=0.25in
}


\usepackage[acronym]{glossaries}
\makeglossaries
\newacronym{ddil}{DDIL}{degraded, denied, intermittent, or limited}
\newacronym{uas}{UAS}{unmanned aerial system}
\newacronym{uav}{UAV}{unmanned aerial vehicle}


\title{Modeling Connectivity-Constrained Dolev-Yao-Controlled Networks for Mobile and Ad Hoc Environments}
\author{Joseph W. Lukefahr}
\date{\today}


\begin{document}

\maketitle

\section{Rationale for Adding Connectivity Contraints to the Dolev-Yao Model}

In the standard Dolev-Yao model, the attacker completely controls message delivery.
Any outgoing message is immediately placed into the attacker's knowledge, and the attacker can replay it for any process at any time. 
Furthermore, in the default model, the attacker can force a process to receive a message at any time. \cite{Dolev_1983_SecurityPublicKey}
In other words, the Dolev-Yao model assumes assumes a perfectly connected network where all processes are always "online". 
The attacker can force message receipt regardless of whether the process should be reachable at that moment.

This is unrealistic for wireless, mobile, ad hoc (e.g., multi-vehicle \gls{uas}), or networks operating in \gls{ddil} environments.
These networks encounter intermittent wireless links, nodes moving in/out of communication range, and network partitions.
Furthermore, network operators may impose deliberate shutdown of radios or limited communication windows (e.g., timed \gls{uav} to satellite transmissions).

We therefore need an explicit notion of online connectivity.
We need a more realistic model where message reception and transmission depend on local network state.

The proposed connectivity-constrained Dolev-Yao model makes connectivity a first-class concept that limits attacker capabilities.
The attacker still controls timing and content of messages (as in standard Dolev–Yao).
However, it cannot violate connectivity assumptions or force message delivery when a process should not be reachable.

\section{On Asynchronicity of this Model}

This model is asynchronous in the sense that there is no global clock or timing assumptions.
Processes can transition between online and offline states at any time, and the attacker can choose when to deliver messages.
However, the model does not include explicit timing constraints or assumptions about message delivery delays.
The attacker can delay message delivery indefinitely, but it cannot force delivery when the recipient is offline.
This allows us to model a wide range of real-world scenarios, including those with intermittent connectivity, without imposing unrealistic timing assumptions.

This asynchronous timing model has a rich history in the literature of security protocol analysis.
Asynchronocity is implied in the original Dolev-Yao model, which does not assume any timing guarantees for message delivery or process execution \cite{Dolev_1983_SecurityPublicKey}.
Since then, several bodies of work have used the Dolev-Yao model and its asynchronous timing model.
Analyses of the Messaging Layer Security continuous group key agreement protocol use this asynchronous timing model \cite{CohnGordon_2018_EndstoEnds, Alwen_2020_Security, Brzuska_2022_Security}.

Asynchronicity has an even richer history in the broader evolution of distributed systems analysis.
It is considered one of the fundamental models of distributed systems, alongside synchronous and partially synchronous models \cite{Lynch_1996_DistributedAlgorithms}.
For example, it founds the Fischer-Lynch-Paterson impossibility result regarding distributed consensus with one faulty process \cite{Fischer_1985_Impossibility}.

\section{Tamarin Ruleset for Connectivity-Constrained Dolev-Yao}

We have four Tamarin rules that model this network.
Two rules model connectivity changes for processes.
Two rules model message transmission/reception in a network where the attacker controls message delivery but is constrained by the connectivity state of the processes.

\paragraph{Connect} This rule models a process transitioning from an offline 
state to an online state. The premise \texttt{Offline(pid)} ensures the rule only fires 
for participants that are currently disconnected. The action fact \texttt{Connect(pid)} 
records that this transition occurred and is useful in proofs about connectivity 
changes or liveness. The conclusion introduces the fact \texttt{Online(pid)}, marking 
that this participant is now able to send and receive messages over the network.

\begin{minipage}{\linewidth}
\lstinputlisting[
    style=lst-default,
    firstnumber=9,
    firstline=9,
    lastline=16
]{main.spthy}
\end{minipage}

\paragraph{Disconnect} This rule models a participant transitioning from being connected to losing 
connectivity. When \texttt{Online(pid)} holds, the rule may fire, producing the action 
fact \texttt{Disconnect(pid)} and updating the participant's state back to 
\texttt{Offline(pid)}. After this transition, the participant should no longer be able 
to send or receive messages unless it later executes a \texttt{Connect} rule.

\begin{minipage}{\linewidth}
\lstinputlisting[
    style=lst-default,
    firstnumber=18,
    firstline=18,
    lastline=25
]{main.spthy}
\end{minipage}

\paragraph{Broadcast Message} This rule models sending a message into the network, 
but only if the sender is online. It requires both \texttt{Online(pid)} and a protocol 
fact \texttt{Message\_To\_Send(pid, m)} to indicate that the node intends to broadcast 
message \texttt{m}. The event \texttt{Broadcast\_Message(pid, m)} marks that the broadcast 
occurred, and \texttt{Out(m)} sends the message into the Dolev–Yao network where the 
attacker gains control of it.

\begin{minipage}{\linewidth}
\lstinputlisting[
    style=lst-default,
    firstnumber=27,
    firstline=27,
    lastline=35
]{main.spthy}
\end{minipage}

\paragraph{Receive Message} This rule models message reception, but with a connectivity constraint.  
The \texttt{In(m)} premise is Tamarin’s standard way of modeling that the attacker has 
chosen to deliver message \texttt{m} to this process at this moment. The presence of 
\texttt{Online(pid)} ensures this delivery can only occur when the node has network 
connectivity. When the rule fires, it generates a trace event 
\texttt{Receive\_Message(pid, m)} and records that the process received the message by 
asserting \texttt{Message\_Received(pid, m)}.

\begin{minipage}{\linewidth}
\lstinputlisting[
    style=lst-default,
    firstnumber=37,
    firstline=37,
    lastline=45
]{main.spthy}
\end{minipage}

\section{Looking Ahead: Process ID as a Function of Signing Key}

The network model described above relies on the concept of process identities (pids) to track which processes are online and to associate messages with specific processes.
In the proposed protocol under development, each process identity could be derived from its signing key (by hashing the corresponding public verification key), specifically as the hash of the signing key:

\begin{equation}
pid = h(pk(sk))
\end{equation}

This design choice has several implications.  
It ensures unique identifiers for each process.  
It also binds the process identity to its fresh cryptographic key material.  
Additionally, it simplifies identity management.

Since signing keys are unique, deriving the process ID from the signing key ensures that each process has a unique identifier. This prevents identity collisions and simplifies authentication.

By tying the process ID to the signing key, we create a strong link between a process's identity and its cryptographic credentials. This can enhance security by making it more difficult for an attacker to impersonate a process without access to the corresponding signing key.

This approach simplifies the management of identities, as there is no need for a separate identity management system or mechanism. The process ID can be easily computed from the signing key.   More importantly, this choice eliminates the need for an entity authentication mechanism.

\subsection{Injectivity in the Signing Key to Process ID Mapping}

We assert that the mapping from signing keys to process IDs is injective, meaning that each signing key corresponds to exactly one process ID.
Specifically, we assert that this function produces the same output only if the input signing keys are the same.
This is the symbolic analog of collision resistance in the underlying hash function.
Collision resistance, a complexity-theoretic notion, ensures that it is hard to find two distinct inputs $x \neq y$ such that $h(x)=h(y)$.
In the symbolic model, we assume that the hash function behaves ideally, meaning that it produces unique outputs for different inputs.
This requires implementations of this proposed protocol to be careful in their choice of hash function, ensuring it has strong collision resistance to uphold the security assumptions of the model.

\printglossaries
\printbibliography

\end{document}